% GENERATED FILE. Edit canonical lessons, then run npm run generate:books.
% canonical-source-hash: fnv1a64:8d2b1a6b1ac383ee
% canonical-lessons: MW-C91-pachhai, MW-C91-somvar, MW-C91-mangalvar, MW-C91-guruvar, MW-C91-sukarvar

\chapter{Afterwards, Monday, Tuesday, Thursday, Friday}
\label{ch:mw-afterwards-monday-tuesday-thursday-friday}
\hlchaptermodality{91}

\begin{chapteropening}
\textbf{By the end of this chapter:} \emph{I can say afterwards, Monday, Tuesday, Thursday and Friday.}

Complete the last lesson of chapter 91: I can say afterwards, Monday, Tuesday, Thursday and Friday.
\end{chapteropening}

\section[\texorpdfstring{\mw{पछै} (pachhai)}{पछै (pachhai)}]{\mw{पछै} (pachhai) — afterwards, later}

\label{lesson:MW-C91-pachhai}



\begin{quote}
Before the new one: say the Marwadi for dawn, then the Marwadi for the day after tomorrow; the day before yesterday.
\end{quote}

\subsection*{You'll want to know: \mw{पछै}}
\textbf{\mw{पछै}} — \emph{pachhai} — \textquotedblleft{}afterwards, later\textquotedblright{}.

\begin{cousinweb}
\textbf{\mw{पछै}} (\emph{pachhai}) is a tadbhava, a word inherited and reshaped by sound change. It comes from \emph{*paśca}, \textquotedblleft{}hinder part\textquotedblright{}, a form CDIAL reconstructs from its descendants. CDIAL lists Old Marwari \emph{pāchaï} under it (CDIAL 7990). CDIAL says this form was drawn from Sanskrit \emph{paścāt}, \textquotedblleft{}behind, afterwards\textquotedblright{}.
\end{cousinweb}

\begin{grammarlens}[title={What you've built}]
One more word for this chapter.
\end{grammarlens}

\subsection*{Your turn}
\begin{itemize}
\raggedright
  \item \emph{Say it:} \emph{pachhai}
  \item \emph{Say it:} \emph{pachhai}, once more
  \item \emph{Say it:} say \emph{pachhai}
  \item \emph{Recall:} say the Marwadi for dawn, then the Marwadi for the day after tomorrow; the day before yesterday, then say \emph{pachhai} again
\end{itemize}

\begin{tcolorbox}[breakable,colback=teal!4,colframe=teal!35!black,title={Before you move on}]
What does \mw{पछै} mean? (\textquotedblleft{}Afterwards, later\textquotedblright{}.) Say it once more.
\end{tcolorbox}
\section[\texorpdfstring{\mw{सोमवार} (somvār)}{सोमवार (somvār)}]{\mw{सोमवार} (somvār) — Monday}

\label{lesson:MW-C91-somvar}



\begin{quote}
Before the new one: say the Marwadi for the day after tomorrow; the day before yesterday, then the Marwadi for afterwards.
\end{quote}

\subsection*{You'll want to know: \mw{सोमवार}}
\textbf{\mw{सोमवार}} — \emph{somvār} — \textquotedblleft{}Monday\textquotedblright{}.

\begin{cousinweb}
\textbf{\mw{सोमवार}} (\emph{somvār}) is a Sanskrit compound, a tatsama: \emph{soma}, \textquotedblleft{}the moon\textquotedblright{} (CDIAL 13608), and \emph{vāra}, \textquotedblleft{}day of the week\textquotedblright{}, whose inherited Old Marwari form is \emph{bāra}, \textquotedblleft{}turn, time\textquotedblright{} (CDIAL 11547).
\end{cousinweb}

\begin{grammarlens}[title={What you've built}]
One more word for this chapter.
\end{grammarlens}

\subsection*{Your turn}
\begin{itemize}
\raggedright
  \item \emph{Say it:} \emph{somvār}
  \item \emph{Say it:} \emph{somvār}, once more
  \item \emph{Say it:} say \emph{somvār}
  \item \emph{Recall:} say the Marwadi for the day after tomorrow; the day before yesterday, then the Marwadi for afterwards, then say \emph{somvār} again
\end{itemize}

\begin{tcolorbox}[breakable,colback=teal!4,colframe=teal!35!black,title={Before you move on}]
What does \mw{सोमवार} mean? (\textquotedblleft{}Monday\textquotedblright{}.) Say it once more.
\end{tcolorbox}
\section[\texorpdfstring{\mw{मंगळवार} (maṅgaḷvār)}{मंगळवार (maṅgaḷvār)}]{\mw{मंगळवार} (maṅgaḷvār) — Tuesday}

\label{lesson:MW-C91-mangalvar}



\begin{quote}
Before the new one: say the Marwadi for afterwards, then the Marwadi for Monday.
\end{quote}

\subsection*{You'll want to know: \mw{मंगळवार}}
\textbf{\mw{मंगळवार}} — \emph{maṅgaḷvār} — \textquotedblleft{}Tuesday\textquotedblright{}.

\begin{cousinweb}
\textbf{\mw{मंगळवार}} (\emph{maṅgaḷvār}) is a Sanskrit compound: \emph{maṅgala}, \textquotedblleft{}the planet Mars\textquotedblright{} (CDIAL 9706), and \emph{vāra}, \textquotedblleft{}day of the week\textquotedblright{} (CDIAL 11547). Marwadi says the \emph{l} curled back.
\end{cousinweb}

\begin{grammarlens}[title={What you've built}]
One more word for this chapter.
\end{grammarlens}

\subsection*{Your turn}
\begin{itemize}
\raggedright
  \item \emph{Say it:} \emph{maṅgaḷvār}
  \item \emph{Say it:} \emph{maṅgaḷvār}, once more
  \item \emph{Say it:} say \emph{maṅgaḷvār}
  \item \emph{Recall:} say the Marwadi for afterwards, then the Marwadi for Monday, then say \emph{maṅgaḷvār} again
\end{itemize}

\begin{tcolorbox}[breakable,colback=teal!4,colframe=teal!35!black,title={Before you move on}]
What does \mw{मंगळवार} mean? (\textquotedblleft{}Tuesday\textquotedblright{}.) Say it once more.
\end{tcolorbox}
\section[\texorpdfstring{\mw{गुरुवार} (guruvār)}{गुरुवार (guruvār)}]{\mw{गुरुवार} (guruvār) — Thursday}

\label{lesson:MW-C91-guruvar}



\begin{quote}
Before the new one: say the Marwadi for Monday, then the Marwadi for Tuesday.
\end{quote}

\subsection*{You'll want to know: \mw{गुरुवार}}
\textbf{\mw{गुरुवार}} — \emph{guruvār} — \textquotedblleft{}Thursday\textquotedblright{}.

\begin{cousinweb}
\textbf{\mw{गुरुवार}} (\emph{guruvār}) is a Sanskrit compound, a tatsama: \emph{guru}, \textquotedblleft{}heavy\textquotedblright{}, also the name of the planet Jupiter (CDIAL 4209), and \emph{vāra}, \textquotedblleft{}day of the week\textquotedblright{} (CDIAL 11547).
\end{cousinweb}

\begin{grammarlens}[title={What you've built}]
One more word for this chapter.
\end{grammarlens}

\subsection*{Your turn}
\begin{itemize}
\raggedright
  \item \emph{Say it:} \emph{guruvār}
  \item \emph{Say it:} \emph{guruvār}, once more
  \item \emph{Say it:} say \emph{guruvār}
  \item \emph{Recall:} say the Marwadi for Monday, then the Marwadi for Tuesday, then say \emph{guruvār} again
\end{itemize}

\begin{tcolorbox}[breakable,colback=teal!4,colframe=teal!35!black,title={Before you move on}]
What does \mw{गुरुवार} mean? (\textquotedblleft{}Thursday\textquotedblright{}.) Say it once more.
\end{tcolorbox}
\section[\texorpdfstring{\mw{सुकरवार} (sukarvār)}{सुकरवार (sukarvār)}]{\mw{सुकरवार} (sukarvār) — Friday}

\label{lesson:MW-C91-sukarvar}



\begin{quote}
Before the new one: say the Marwadi for Tuesday, then the Marwadi for Thursday.
\end{quote}

\subsection*{You'll want to know: \mw{सुकरवार}}
\textbf{\mw{सुकरवार}} — \emph{sukarvār} — \textquotedblleft{}Friday\textquotedblright{}.

\begin{cousinweb}
\textbf{\mw{सुकरवार}} (\emph{sukarvār}) is a Sanskrit compound, lightly reshaped, a semi-tatsama: \emph{śukra}, \textquotedblleft{}bright\textquotedblright{}, also the planet Venus (CDIAL 12506), and \emph{vāra}, \textquotedblleft{}day of the week\textquotedblright{} (CDIAL 11547). Marwadi says the \emph{ś} as \emph{s} and splits the \emph{kr} with a vowel.
\end{cousinweb}

\begin{grammarlens}[title={What you've built}]
One more word for this chapter.
\end{grammarlens}

\subsection*{Your turn}
\begin{itemize}
\raggedright
  \item \emph{Say it:} \emph{sukarvār}
  \item \emph{Say it:} \emph{sukarvār}, once more
  \item \emph{Say it:} say \emph{sukarvār}
  \item \emph{Recall:} say the Marwadi for Tuesday, then the Marwadi for Thursday, then say \emph{sukarvār} again
\end{itemize}

\begin{tcolorbox}[breakable,colback=teal!4,colframe=teal!35!black,title={Before you move on}]
What does \mw{सुकरवार} mean? (\textquotedblleft{}Friday\textquotedblright{}.) Say it once more.
\end{tcolorbox}
